
2.3.
Nous aurons besoin d'appliquer le foncteur \(\Projz\) à des \(\underline{B}\)-Modules gradués quelconques, et non seulement à des \(\underline{B}\)-Modules gradués quasi-cohérents comme dans la construction de EGA II 3.2. Aussi allons nous indiquer rapidement comment cette construction peut s'effectuer dans le cas général.

2.3.1. Soient d'abord \(S\) un schéma, \(\underline{A}\) une \(\underline{O}_S\)-algèbre quasi-cohérente, \(Y=\Spec(\underline{A})\), et \(f:Y\to S\) le morphisme structural. De l'isomorphisme~:
\[
  \underline{A}\xrightarrow{\sim} f_*(\underline{O}_Y)
\]
on déduit un homomorphisme canonique~:
\begin{equation}\tag{2.3.1.1}
  f^{-1}(\underline{A})\longrightarrow \underline{O}_Y
\end{equation}
où \(f^{-1}\) désigne l'image inverse au sens des faisceaux d'ensembles; en d'autres termes, on a un morphisme d'espaces annelés~:
\begin{equation}\tag{2.3.1.2}
  \bar f:(Y,\underline{O}_Y)\longrightarrow (S,\underline{A})
\end{equation}
On en déduit un foncteur «image inverse» \(\bar f^{\,*}\) de la catégorie des \(\underline{A}\)-Modules dans la catégorie des \(\underline{O}_Y\)-Modules, adjoint à gauche du foncteur «image directe» \(\bar f_*\) de la catégorie des \(\underline{O}_Y\)-Modules dans la catégorie des \(f_*(\underline{O}_Y)\)-Modules. On a donc~:
\begin{equation}\tag{2.3.1.3}
  \bar f^{\,*}(\underline{M})=f^{-1}(\underline{M})\otimes_{f^{-1}(\underline{A})}\underline{O}_Y
\end{equation}
pour tout \(\underline{A}\)-Module \(\underline{M}\).

Comme l'homomorphisme (2.3.1.1) est plat (car dans chaque fibre il correspond à une localisation), le foncteur \(\bar f^{\,*}\) ainsi défini est un foncteur exact (cf [4], p. 164).

Si \(\underline{M}\) est un \(\underline{A}\)-Module quasi-cohérent, cette construction donne le Module \(\widetilde{\underline{M}}\) défini dans EGA II 1.4.3. En effet, l'isomorphisme~:
